\begin{textsample}{POS Dim 3 – global\_north\_2024\_12 – Score 15.00 – global\_north\_2024\_12/118439300.txt}  \label{ex:f3_pos_237}
Oxford and other top \textbf{universities} must increase efforts to combat antisemitism, a former Oxford professor has said.
Sir Vernon Bogdanor, a former professor of politics and government at the \textbf{University} of Oxford and now a professor of government at King ' s \textbf{College} London, said \textbf{universities} failing to tackle antisemitism and discrimination in violation of the Equality Act 2010 should face cuts to government funding.
The letter claimed that the Oxford Union debate on Israel, which took place last Thursday and passed a motion labelling Israel as an"apartheid state responsible for genocide", violated the law by subjecting Jewish \textbf{students} to antisemitism and intimidation.
Bogdanor, 81, said :"It wasn't a debate in any sense.
It was a kangaroo court with comments which were \textbf{antisemitic} and were against the law.
Advertisement"It is also fairly clear that there was so much intimidation that several \textbf{students} were rather frightened to attend.
So, it wasn't really a proper debate of the sort that the Oxford Union ought to have."Vernon Bogdanor said both the \textbf{university} and union have not done enough to tackle antisemitism ALAMY Discrimination based on religion or belief is illegal under the Equality Act 2010, while the Crime and Disorder Act 1998 defines hostility or incitement of racial hatred as a hate crime.
At the debate, one speaker reportedly praised an attack on Israel as"heroism", prompting criticism of the event ' s"inflammatory rhetoric, aggressive behaviour, and intimidation".
Jonathan Sacerdoti, the opposing speaker and son of a Holocaust survivor, accused the Union of"disgracing itself"by amplifying"bigotry, hatred, and mob rule".
Bogdanor added :"There are no grounds for stopping a debate on the Gaza war and Israeli policies in the Union.
There should be such a debate and \textbf{students} are entitled to feel strongly on both sides of the issue.
But as I said, there was freedom of \textbf{speech} for those who were hostile to Israeli policies and not for those who wished to defend them and that ' s not a debate...
It casts a light on the attitudes of \textbf{students}."Advertisement While the \textbf{university} maintains a neutral stance on international conflicts, Bogdanor argues that both the \textbf{institution} and the Oxford Union -- although constitutionally separate -- have failed to do enough to tackle discrimination.
He said :"I don't think they ' ve done enough.
They produced a rather generalised document about racism but I think they have not been proactive enough in not letting criticism about Israel generate antisemitism.
It ought to be punished where it is there and anyone who stops freedom of \textbf{speech} -- which is what is happening in the Union -- should be sent down."Bogdanor said Oxford should publish a formal document on antisemitism, referencing the internationally accepted definition and explaining its unacceptability.
He added :"I don't think Oxford is untypical sadly.
I think it ' s true of many \textbf{universities}, particularly elite \textbf{universities}.
As Andrew Neil said, ' The more elite the British \textbf{university}, the more stupid the \textbf{students} ', and I think the same unfortunately is true of elite \textbf{universities} in America."The conflict in Gaza has also prompted unrest at Cambridge, where \textbf{graduation} \textbf{ceremonies} for hundreds of \textbf{students} were disrupted by \textbf{pro-Palestinian} \textbf{protests}.
Advertisement The Cambridge for Palestine movement received a letter from the \textbf{university} warning of potential legal action and exclusion.
The \textbf{university} acknowledged the right to \textbf{peaceful} \textbf{protest} but condemned the disruption caused by the occupations.
In May, ministers called on \textbf{universities} to crack down on \textbf{antisemitic} abuse on \textbf{campuses} and to ensure \textbf{protests} do not disrupt \textbf{university} life.
Bogdanor thinks Sir Keir Starmer ' s government should do the same.
He said :"Rishi Sunak reminded \textbf{universities} of their responsibilities and I think Keir Starmer should do the same.
It ' s interesting that Donald Trump has told \textbf{universities} that if they don't deal with this issue, they ' ll lose \textbf{federal} funding and I don't think Keir Starmer should appear to be less diligent in dealing with racism than Donald Trump."A \textbf{university} spokesman said :"The Oxford Union is not part of the \textbf{university} and must speak for itself about its events, which are neither regulated nor recognised by the \textbf{university}."The \textbf{University} rejects and condemns antisemitism and makes this clear in all its published policies, adopting the IHRA definition of antisemitism.
There is no place for unlawful discrimination of any kind directed towards any faith, race, nationality or ethnic group at Oxford.
We are committed to ensuring that all \textbf{students} feel safe and included, and investigate all formal complaints made about \textbf{harassment} or discrimination at the \textbf{university}.
Advertisement"\textbf{University} leaders and \textbf{officers} have been working with Jewish \textbf{students} and staff to strengthen our response to \textbf{harassment} and discrimination.
Already this \textbf{academic} year, we have delivered new training to \textbf{student} leaders, freshers and staff that has included attention to antisemitism.
We have also established a task and finish group to consider \textbf{student} experiences relating to racial and religious discrimination and recommend ways in which we can improve our response."We have made very clear our expectation that all members of our \textbf{university} must conduct themselves with respect, civility and empathy."Universities UK declined to comment, saying it was not"able to comment on individual \textbf{institutions}".

% matched lemmas: academic, antisemitic, campus, ceremony, college, federal, graduation, harassment, institution, officer, peaceful, pro-palestinian, protest, speech, student, university
\end{textsample}
